% =============================================================================
%  VERSION INSTITUT DES ACTUAIRES A 150 PAGES (octobre 2026)
%
%  Troisieme fichier maitre, a cote de main_ensae.tex (version SFE, 101 pages) et de
%  main_institut.tex (version deposee le 30 septembre, 195 pages, FIGEE). Il lit
%  chapitres_inst/, copie de chapitres_inst/ sur laquelle sont ajoutes les contenus qui
%  font passer le document de 100 a 150 pages sans perte de qualite. Il ne modifie
%  aucun fichier lu par les deux autres versions.
%
%  Differences avec la version ENSAE : couverture officielle de l'Institut
%  (page_de_garde), paragraphe d'usage de l'IA a la fin du resume (§ 5.1.c des
%  Recommandations Jury), notes de synthese en tete, pas de chapitre du stage,
%  annexes F (pistes) et G (inventaire de l'IA) exigees par l'Institut.
%  Echeance cote Institut : janvier 2027 (soutenance en mars 2027).
% =============================================================================
\documentclass[12pt,a4paper]{report}

% FORMAT IMPOSE PAR L'ECOLE : Times New Roman 12 et interligne 1,5. La commande ci-dessous
% active, dans le preambule partage, la branche Times et 1,5 ; les autres versions ne la
% definissent pas et gardent Palatino 11. L'ecole a leve la limite de trente pages le
% 16 septembre 2026, pas la police ni l'interligne.
\newcommand{\formatensae}{}

% Meme preambule que la version 2, qui est la reference de style du projet.
\input{preambule_v2}

% Laisser le haut de page libre autour du logo du fond Nexialog.
% Le numero de page reste en haut a droite.
\fancyhead{}
\fancyhead[R]{\small\thepage}
\renewcommand{\headrulewidth}{0pt}

% Fond de page Nexialog : fichier present dans le projet Overleaf.
\usepackage{eso-pic}
\AddToShipoutPictureBG{%
  \AtPageLowerLeft{%
    \includegraphics[width=\paperwidth,height=\paperheight]{RequiredFiles/fond_latex_2.pdf}%
  }%
}

% Espaces de lecture propres à cette version : police 12 et interligne 1,5 conservés.
% Les définitions restent près des équations ; les blancs autour des sous-titres sont resserrés.
\titlespacing*{\section}{0pt}{14pt plus 2pt minus 2pt}{7pt plus 1pt}
\titlespacing*{\subsection}{0pt}{11pt plus 2pt minus 1pt}{5pt plus 1pt}

% Chemins des images depuis la racine du projet Overleaf.
% Le dossier figures_memoire/ provient du ZIP correctif des images.
% Les derniers chemins permettent aussi d'utiliser l'organisation du ZIP initial.
\graphicspath{%
  {./}%
  {figures_memoire/}%
  {figures_pedagogie/}%
  {figures_externes/}%
  {logos/}%
  {RequiredFiles/}%
  {recherche/vasicek_lab/figures_inst/}%
  {recherche/vasicek_lab/figures/}%
  {recherche/cascade_qualitative/figures/}%
  {resultats/outputs/}%
}

% VERSION ENSAE SEULE (reduction SFE, phase 3). Deux reglages qui ne touchent pas le preambule
% partage avec la version Institut :
%  - la bibliographie entre au sommaire, que les consignes MS exigent et ou natbib ne l'inscrit
%    pas ;
%  - un renvoi vers un contenu parti au supplement passe par une ETIQUETTE RELAIS, posee dans la
%    notice des annexes : le \ref du texte reste intact et imprime la place du contenu dans le
%    supplement (S1, S1.2...), le lien menant a la notice qui dit ou le trouver.
\renewcommand{\bibsection}{\chapter*{\bibname}\phantomsection
  \addcontentsline{toc}{chapter}{\bibname}\markboth{\bibname}{}}
\makeatletter
\newcommand{\etiquettesupplement}[2]{\phantomsection\def\@currentlabel{#2}\label{#1}}
\makeatother

\begin{document}

% COUVERTURE : le gabarit officiel de l'Institut des actuaires (page_de_garde.tex),
\newcommand*{\pgdatesoutenance}{mars 2027}
% Titre : celui de la version ENSAE, valeurs par défaut du gabarit (décision de Kélian, 29 octobre 2026).
\input{page_de_garde}

% PIECES LIMINAIRES : remerciements, resume, abstract, sommaire, glossaire. Copies de
% remerciements.tex, 01_resume.tex et 00_glossaire.tex, a reecrire en phase 4.
\input{chapitres_inst/remerciements}
\input{chapitres_inst/resume}      % Résumé, abstract, sommaire
\input{chapitres_inst/glossaire}
\input{chapitres_inst/notes_synthese}   % notes de synthese FR et EN, en tete

% CORPS, SELON LE PLAN CIBLE DE LA MISSION. Plus de parties : leurs six etiquettes part:...
% sont reposees sur le chapitre qui ouvre desormais leur contenu.
\input{chapitres_inst/00_introduction}             % ex-1.5.3
\input{chapitres_inst/01_contexte}                 % ex-chapitres 1 et 2
\input{chapitres_inst/02_revue}                    % ex-3
\input{chapitres_inst/03_donnees}                  % ex-4
\input{chapitres_inst/04_frequence_severite}       % ex-5
\input{chapitres_inst/05_cascade}                  % ex-6
\input{chapitres_inst/06_multietats}               % ex-7
\input{chapitres_inst/07_identifiabilite}          % ex-8
\input{chapitres_inst/08_identification_partielle} % ex-9
\input{chapitres_inst/09_resultats}                % ex-10
\input{chapitres_inst/10_robustesse}               % ex-11 et 12
\input{chapitres_inst/conclusion}                  % ex-13

% BIBLIOGRAPHIE AVANT LES ANNEXES : ordre des consignes MS (sommaire, introduction,
% developpements, conclusion, bibliographie, annexes).
\bibliographystyle{plainnat-fr}
% Interligne simple et entrees resserrees dans la bibliographie (1er octobre 2026, demande de
% Kelian : faire la place des figures remontees du supplement). La police reste le Times 12 des
% consignes MS ; l'interligne 1,5 vaut pour le texte courant, comme dans les tableaux.
% Corps \footnotesize depuis le 28 octobre 2026 : rend la page prise par les notes de synthese a trois pages.
{\NoAutoSpacing\singlespacing\footnotesize\setlength{\bibsep}{0pt plus 0.5pt}\bibliography{references_inst}}

% =====================================================================
\appendix
% Plus de page de partie : une ligne « Annexes » au sommaire, en tete de la notice.
\clearpage
\phantomsection
\addcontentsline{toc}{part}{Annexes}

% -----------------------------------------------------------------------------
% NOTICE D'ORGANISATION DES ANNEXES. Les annexes B (demonstrations du socle EVT),
% D (pieces justificatives) et E (complements aux chapitres du corps) ont ete
% deportees le 18 septembre 2026 vers materiel_complementaire/ et ne sont plus
% appelees ; les renvois vers elles passent par la macro \matcomp. Depuis la phase 3
% de la reduction, trois autres contenus sont partis dans supplement/ (S1 a S3) ; les
% quatre etiquettes que le corps vise encore y sont relayees ci-dessous.
% -----------------------------------------------------------------------------
\chapter*{Organisation des annexes et matériel complémentaire}
\addcontentsline{toc}{chapter}{Organisation des annexes et matériel complémentaire}
\markboth{Annexes}{}

Ce mémoire garde en annexe les démonstrations du modèle de cascade (A), celles de la théorie des
valeurs extrêmes avec les détails de validation (B), les adaptations par pilier et le cas d'usage (C), le
protocole d'élicitation et les relectures de praticiens (D), la table des paramètres et des notations
(E), les pistes pour de futurs mémoires (F) et l'inventaire de l'usage de l'intelligence artificielle
générative (G).

Le \emph{matériel complémentaire}, qui réunit les pièces justificatives, les compléments de chapitres et
les diagnostics secondaires, est déposé de façon pérenne sur Zenodo, sous l'identifiant
\texttt{10.5281/zenodo.22874524} :
\url{https://doi.org/10.5281/zenodo.22874524}

Le renvoi numéroté S4.6, au chapitre des données, conduit au supplément de la version rédigée pour
l'ENSAE, déposé au même identifiant dans le fichier \texttt{supplement\_rapport\_ENSAE.pdf}.
% ETIQUETTES RELAIS : fichier ecrit par reduction/relais_supplement.py, a relancer apres tout
% deplacement vers le supplement.
\input{chapitres_inst/relais}

% =====================================================================
% Annexe A revenue du supplement le 1er octobre 2026 tard le soir (« enrichissons le memoire ») :
% c'est l'annexe A du plan cible. Le supplement passe a S1-S4.
\input{chapitres_inst/annexe_a_demonstrations}        % A, ex-A (le supplement l'appelait S5)
\input{chapitres_inst/annexe_evt}                     % B, demonstrations et complements sur les valeurs extremes (ex-materiel complementaire)
\input{chapitres_inst/annexe_adaptations}             % B, adaptations par pilier et cas d'usage (ex-S1)
\input{chapitres_inst/annexe_elicitation}              % C, protocole d'elicitation et relectures de praticiens (ex-S4)
\input{chapitres_inst/annexe_b_parametres_notations}  % D, ex-D et ex-E
\input{chapitres_inst/annexe_pistes}                  % D, exigee par l'Institut
\input{chapitres_inst/annexe_d_ia}                    % E, inventaire de l'IA (exigee)


\end{document}
